\section{Refleksi, Portofolio, dan Tindak Lanjut}
\label{sec:refleksi-tindak-lanjut}

\subsection{Penilaian Diri}

Penilaian diri adalah bagian resmi penilaian. Mahasiswa menilai
pencapaian CPMK-nya sendiri dengan menunjuk bukti. Klaim tanpa bukti
dianggap belum tercapai.

\begin{table}[H]
    \centering
    \caption{Lembar penilaian diri berbasis bukti.}
    \label{tab:self-assessment}
    \small
    \begin{tabular}{@{}p{3.5cm}p{3.25cm}p{3.25cm}p{3.4cm}@{}}
        \toprule
        \textbf{CPMK} & \textbf{Tingkat yang diklaim} &
        \textbf{Bukti} & \textbf{Rencana perbaikan} \\
        \midrule
        \CPMK{01} & belum / cukup / baik / sangat baik &
        Tugas, kuis, UTS & \\
        \addlinespace
        \CPMK{02} & belum / cukup / baik / sangat baik &
        Tugas, praktikum, UTS & \\
        \addlinespace
        \CPMK{03} & belum / cukup / baik / sangat baik &
        Tugas, praktikum, proyek & \\
        \addlinespace
        \CPMK{04} & belum / cukup / baik / sangat baik &
        Laporan \eng{debugging} & \\
        \bottomrule
    \end{tabular}
\end{table}

\subsection{Penilaian Antarteman}

Untuk pekerjaan berpasangan, setiap mahasiswa menilai pasangannya
pada tiga aspek dengan skala 1--4:

\begin{enumerate}
    \item kontribusi teknis terhadap kode;
    \item kualitas kolaborasi;
    \item keandalan memenuhi kesepakatan.
\end{enumerate}

Nilai antarteman tidak dipakai untuk menjatuhkan pasangan. Selisih
besar antara penilaian diri dan penilaian antarteman menjadi bahan
diskusi.

\subsection{Jurnal Refleksi Akhir}

Tulis jurnal refleksi akhir yang menjawab pertanyaan berikut.

\begin{enumerate}
    \item Kompetensi apa yang paling berkembang? Tunjukkan buktinya.
    \item Kompetensi apa yang masih lemah? Apa rencana perbaikannya?
    \item Bagian mana dari proses belajar yang paling tidak efisien,
          dan bagaimana Anda akan mengubahnya?
    \item Setelah memahami level instruksi, bagaimana pandangan Anda
          terhadap program yang Anda tulis dalam bahasa tingkat
          tinggi?
\end{enumerate}

\subsection{Tindak Lanjut dalam Lingkup 8086}

Mata kuliah ini adalah pintu masuk. Mahasiswa dapat memperdalam
materi yang sudah tersedia pada dua buku sumber.

\begin{table}[H]
    \centering
    \caption{Jalur tindak lanjut dalam lingkup 8086.}
    \label{tab:tindak-lanjut}
    \small
    \begin{tabular}{@{}p{3.4cm}p{5.4cm}p{4.8cm}@{}}
        \toprule
        \textbf{Minat} & \textbf{Topik lanjutan} & \textbf{Sumber} \\
        \midrule
        Layanan sistem & Layanan BIOS/DOS lain, operasi layar dan
        \eng{keyboard} & [1] Bagian 18; [2] BAB XVIII \\
        \addlinespace
        Program modular & Prosedur berparameter, macro, dan
        pengelolaan stack & [1] Bagian 12--13; [2] BAB XIV, XVI--XVII \\
        \addlinespace
        Pengolahan data & Array, string, operasi file, dan pencetakan
        bilangan & [1] Bagian 10, 14.1, 15--16; [2] BAB XIX--XXIII \\
        \addlinespace
        \eng{Debugging} & Penelusuran lanjutan dengan DEBUG dan Turbo
        Debugger & [1] Bagian 4.3; [2] BAB XXVII \\
        \bottomrule
    \end{tabular}
\end{table}

\begin{infobox}
\vspace{1.1em}
Bila hanya satu hal yang diingat dari buku ini, ingatlah bahwa
perilaku program selalu dapat ditelusuri: ada urutan instruksi yang
dapat diikuti, register yang dapat diperiksa, dan memori yang dapat
dibaca. Kemampuan menelusuri inilah yang membedakan programmer yang
dapat memperbaiki masalah dari programmer yang hanya menebak.
\end{infobox}

\subsection{Penutup}

Buku ini dimulai dari pertanyaan bagaimana komputer bekerja pada
level instruksi. Sepanjang bab, pertanyaan itu dijawab setahap demi
setahap: dari bilangan dan register, melalui instruksi dan memori,
sampai program yang mengolah array, string, dan bilangan. Bila pada
akhirnya mahasiswa mampu menulis program assembly, menelusuri
jalannya instruksi, dan memperbaiki kesalahan berdasarkan bukti,
maka tujuan buku ini telah tercapai.
